%---------------------------------------------------------------------
\section{Introduction}

Hangboarding involves isometric dead-hangs from edges, pockets, or slopers to overload the finger flexors beyond what routine climbing provides. MFS and forearm endurance are among the strongest measurable predictors of climbing performance in trained cohorts: Pearson correlations of $r = 0.70$--$0.85$ between MFS (or grip endurance) and redpoint grade are consistently reported \cite{watts2004,macleod2007,fryer2016}. General upper-body resistance training does not produce equivalent finger-flexor adaptations, establishing hangboarding as a distinct training stimulus. The first hangboard RCTs appeared circa 2012 \cite{lopezrivera2012,medernach2015}; the current literature comprises ten intervention trials and one large retrospective cohort \cite{lopezrivera2012,medernach2015,lopezrivera2019,levernier2019,mundry2021,devise2022,hermans2022,dindorf2024,langer2024,ruizlopez2025,gilmore2024}, with intervention $n = 11$--54 and cohort $n = 527$, durations of 4--10 weeks. A systematic review and meta-analysis by Stien et al.\ \cite{stien2023} synthesised five of these trials ($n = 110$). Universal limitations include lack of allocation concealment, male-advanced-only samples, and outcomes measured on the hangboard rather than on-wall performance.
